\section{Lesson 08 --- 16/10/2026}
Lesson missed, I did it by myself using the slides.

So far we studied unconstrained problems, where a local minimum must be a stationary point. Now we move to problems with constraints: the condition $\nabla f(x^*) = 0$ is no longer necessary, and we need optimality conditions that also take the constraints into account.

\subsection{First-order optimality conditions for constrained problems}

Consider the constrained optimization problem
\[
    \min f(x), \quad x \in X \coloneqq \{x \in \R^n : g_i(x) \le 0,\ i = 1, \dots, m,\ h_j(x) = 0,\ j = 1, \dots, p\}, \tag{P}
\]
where $f$, $g_i$ and $h_j$ are continuously differentiable for any $i, j$. We write $g = (g_1, \dots, g_m)$ and $h = (h_1, \dots, h_p)$, so the constraints are $g(x) \le 0$ and $h(x) = 0$.

To describe the feasible set around a point $x^*$ we use two cones: one built from the geometry of $X$ and one built from the gradients of the constraints.

\begin{definition}[Tangent cone]
    The \g{tangent cone} at $x^* \in X$ is
    \[
        T_X(x^*) = \left\{ d \in \R^n : \exists\, \{z_k\} \subset X,\ \exists\, \{t_k\} > 0,\ z_k \to x^*,\ t_k \to 0,\ \lim_{k \to \infty} \frac{z_k - x^*}{t_k} = d \right\}.
    \]
\end{definition}

So $T_X(x^*)$ contains the directions along which we can approach $x^*$ while staying in $X$.

\begin{definition}[Active constraints]
    The set of \g{active constraints} at $x^* \in X$ is
    \[
        A(x^*) = \{i : g_i(x^*) = 0\},
    \]
    i.e., the inequality constraints satisfied with equality at $x^*$.
\end{definition}

\begin{definition}[First-order feasible direction cone]
    The \g{first-order feasible direction cone} at $x^* \in X$ is
    \[
        D(x^*) = \left\{ d \in \R^n :
        \begin{aligned}
            & d^T\nabla g_i(x^*) \le 0 && \forall\, i \in A(x^*),\\
            & d^T\nabla h_j(x^*) = 0 && \forall\, j = 1, \dots, p
        \end{aligned}
        \right\}.
    \]
\end{definition}

$D(x^*)$ is the linearization of the feasible set at $x^*$: only the active constraints matter, because the inactive ones remain satisfied in a neighborhood of $x^*$.

\subsection{Abadie constraint qualification (ACQ)}

The tangent cone is the "true" local description of $X$, but it is hard to compute. The cone $D(x^*)$ is easy to compute, since it only needs gradients. We would like the two cones to coincide.

\begin{definition}[Abadie constraint qualification]
    We say that the \g{Abadie constraint qualification} (ACQ) holds at a point $x^* \in X$ if
    \[
        T_X(x^*) = D(x^*).
    \]
\end{definition}

Checking ACQ directly from the definition is not easy either, so the following theorem gives some sufficient conditions that are simple to verify.

\begin{theorem}[Sufficient conditions for ACQ]
    \begin{enumerate}[label=\alph*)]
        \item \g{Affine constraints}: if $g_i$ and $h_j$ are affine for all $i = 1, \dots, m$ and $j = 1, \dots, p$, then ACQ holds at any $x \in X$.
        \item \g{Slater condition} for convex problems: if $g_i$ are convex for all $i = 1, \dots, m$, $h_j$ are affine for all $j = 1, \dots, p$ and there exists $\bar{x} \in X$ s.t.\ $g(\bar{x}) < 0$ and $h(\bar{x}) = 0$, then ACQ holds at any $x \in X$.
        \item \g{Linear independence} of the gradients of active constraints: if $x^* \in X$ and the vectors
        \[
            \nabla g_i(x^*) \text{ for } i \in A(x^*), \qquad \nabla h_j(x^*) \text{ for } j = 1, \dots, p
        \]
        are linearly independent, then ACQ holds at $x^*$.
    \end{enumerate}
\end{theorem}

\subsection{Karush-Kuhn-Tucker theorem}

The first step is a geometric necessary condition: at a local minimum, no direction that stays in $X$ can be a descent direction.

\begin{lemma}
    If $x^*$ is a local minimum of (P), then
    \[
        d^T\nabla f(x^*) \ge 0 \qquad \forall\, d \in T_X(x^*).
    \]
\end{lemma}

This condition uses $T_X(x^*)$, which is hard to handle. If ACQ holds, we can replace $T_X(x^*)$ with $D(x^*)$, and this leads to an algebraic system in terms of the gradients.

\begin{theorem}[Karush-Kuhn-Tucker]
    \label{thm:kkt}
    If $x^*$ is a local minimum and ACQ holds at $x^*$, then there exist $\lambda^* \in \R^m$ and $\mu^* \in \R^p$ s.t.\ $(x^*, \lambda^*, \mu^*)$ satisfies the \g{KKT system}:
    \[
        \begin{cases}
            \nabla f(x^*) + \sum_{i=1}^{m} \lambda_i^* \nabla g_i(x^*) + \sum_{j=1}^{p} \mu_j^* \nabla h_j(x^*) = 0,\\
            \lambda_i^* g_i(x^*) = 0 \quad \forall\, i = 1, \dots, m,\\
            \lambda^* \ge 0,\\
            g(x^*) \le 0,\\
            h(x^*) = 0.
        \end{cases}
    \]
\end{theorem}

The vectors $\lambda^*$ and $\mu^*$ are called \g{KKT multipliers}. The KKT system can be written in a compact form using the following function.

\begin{definition}[Lagrangian function]
    The \g{Lagrangian function} $L : \R^n \times \R^m \times \R^p \to \R$ is
    \[
        L(x, \lambda, \mu) \coloneqq f(x) + \sum_{i=1}^{m} \lambda_i g_i(x) + \sum_{j=1}^{p} \mu_j h_j(x).
    \]
\end{definition}

Then the KKT system can be formulated as:
\[
    \begin{cases}
        \nabla_x L(x, \lambda, \mu) = 0,\\
        \lambda_i g_i(x) = 0, \quad i = 1, \dots, m,\\
        \lambda \ge 0,\ h(x) = 0,\ g(x) \le 0.
    \end{cases}
\]

\begin{remark}
    Since $\lambda \ge 0$ and $g(x) \le 0$, every term $\lambda_i g_i(x)$ is $\le 0$. So the condition $\lambda_i g_i(x) = 0$ for $i = 1, \dots, m$ is equivalent to $\lambda^T g(x) = 0$, or also $\langle \lambda, g(x) \rangle = 0$ (\g{complementarity condition}).
\end{remark}

\begin{example}
    \label{ex:kkt-circle}
    Find the minimum points of the function $f(x_1, x_2) = 2x_1 + x_2$ on the set
    \[
        X = \{(x_1, x_2) \in \R^2 : x_1^2 + x_2^2 \le 4\}.
    \]
    Note that $X$ is compact, so from the Weierstrass theorem $f$ admits a global maximum and a global minimum on $X$.

    The Lagrangian function is
    \[
        L(x_1, x_2, \lambda) = 2x_1 + x_2 + \lambda(x_1^2 + x_2^2 - 4).
    \]
    The KKT system is given by:
    \[
        \begin{cases}
            2 + 2\lambda x_1 = 0,\\
            1 + 2\lambda x_2 = 0,\\
            \lambda(x_1^2 + x_2^2 - 4) = 0,\\
            x_1^2 + x_2^2 \le 4,\\
            \lambda \ge 0.
        \end{cases}
    \]
    Note that for $\lambda = 0$ the system is impossible (the first equation becomes $2 = 0$). So $\lambda > 0$, the constraint is active, and we are led to solve the system:
    \[
        \begin{cases}
            x_1 = -\dfrac{1}{\lambda},\\[2mm]
            x_2 = -\dfrac{1}{2\lambda},\\[2mm]
            x_1^2 + x_2^2 = 4,\\
            \lambda \ge 0.
        \end{cases}
    \]
    Then
    \[
        \frac{1}{\lambda^2} + \frac{1}{4\lambda^2} = 4,
    \]
    from which
    \[
        16\lambda^2 = 5, \qquad \lambda = \pm \frac{\sqrt{5}}{4}.
    \]
    Since $\lambda \ge 0$, we only keep $\lambda = \frac{\sqrt{5}}{4}$:
    \[
        x_1 = -\frac{4}{\sqrt{5}} = -\frac{4\sqrt{5}}{5}, \qquad x_2 = -\frac{2}{\sqrt{5}} = -\frac{2\sqrt{5}}{5}.
    \]
    The global minimum exists and, as shown in the remark below, ACQ holds at every point of $X$, so by \cref{thm:kkt} the global minimum must solve the KKT system. Since the KKT system has only one solution, it follows that $\bar{x} = \left(-\frac{4\sqrt{5}}{5}, -\frac{2\sqrt{5}}{5}\right)$ is the global minimum point.
\end{example}

\begin{remark}
    In the previous example ACQ is fulfilled for every $x \in X$. Indeed, there is only the constraint $g(x) = x_1^2 + x_2^2 - 4 \le 0$, with $\nabla g(x_1, x_2) \neq (0, 0)$ for every $x \in X$ s.t.\ $x_1^2 + x_2^2 - 4 = 0$. So condition c) of the sufficient conditions for ACQ holds.
\end{remark}

What about the maximum points of $f$ on $X$? It is enough to set $\lambda \le 0$ in the KKT system. In fact, to find the maxima of $f$ we look for the minima of $-f$. The KKT system for the problem
\[
    \min\, (-f(x)), \quad x \in X
\]
is
\[
    \begin{cases}
        -2 + 2\lambda x_1 = 0,\\
        -1 + 2\lambda x_2 = 0,\\
        \lambda(x_1^2 + x_2^2 - 4) = 0,\\
        x_1^2 + x_2^2 \le 4,\\
        \lambda \ge 0,
    \end{cases}
\]
which is equivalent to setting $\lambda \le 0$ in the original one. Choosing $\lambda = -\frac{\sqrt{5}}{4}$ in the original system we obtain
\[
    x_1 = \frac{4}{\sqrt{5}} = \frac{4\sqrt{5}}{5}, \qquad x_2 = \frac{2}{\sqrt{5}} = \frac{2\sqrt{5}}{5},
\]
so that $\hat{x} = \left(\frac{4\sqrt{5}}{5}, \frac{2\sqrt{5}}{5}\right)$ is the global maximum point of $f$ on $X$.

\subsubsection{Limits of the KKT theorem}

The KKT theorem has two limits: it needs ACQ, and it only gives necessary conditions.

\begin{remark}
    The ACQ assumption is crucial in the KKT theorem.
\end{remark}

\begin{example}
    Consider
    \[
        \begin{cases}
            \min\, x_1 + x_2\\
            (x_1 - 1)^2 + (x_2 - 1)^2 - 1 \le 0\\
            x_2 \le 0
        \end{cases}
    \]
    The feasible set is the single point $x^* = (1, 0)$, which is therefore the global optimum.

    $T_X(x^*) = \{0\}$ and $D(x^*) = \{d \in \R^2 : d_2 = 0\}$, hence ACQ does not hold at $x^*$.

    $\nabla g_1(x^*) = (0, -2)$, $\nabla g_2(x^*) = (0, 1)$, $\nabla f(x^*) = (1, 1)$, hence there is no $\lambda^*$ s.t.\ $(x^*, \lambda^*)$ solves the KKT system: the first component of $\nabla_x L$ is always $1 \neq 0$.
\end{example}

The KKT theorem gives necessary optimality conditions, but not sufficient ones.

\begin{example}
    Consider
    \[
        \begin{cases}
            \min\, x_1^2 - x_2^2\\
            x_1^2 + x_2^2 - 2 \le 0
        \end{cases}
    \]
    $x^* = (0, 0)$, $\lambda^* = 0$ solves the KKT system, but $x^*$ is not a local optimum.
\end{example}

For convex problems, instead, the KKT conditions are also sufficient.

\begin{theorem}
    \label{thm:kkt-convex}
    If the optimization problem is convex and $(x^*, \lambda^*, \mu^*)$ solves the KKT system, then $x^*$ is a global optimum.
\end{theorem}

Recall that (P) is convex if $f$ and $g$ are convex and $h$ is affine.

\s{Exercise 4.1.} Prove \cref{thm:kkt-convex}.
